\section{A Business Is a Cooperation Machine}

\subsection{Coordination and the Boundary of the Firm}

Founders tend to receive disproportionate public credit, partly because individuals make stories easier to tell. A company, however, is a coordinated system of people and assets.

The engineer who invents something, the salesperson who earns a customer's trust, the operator who prevents failures, the accountant who keeps the firm solvent, and the customer who offers frank criticism each contribute something different. The founder's role is not to be better at every task. It is to bring capabilities together around a valuable objective and establish a structure in which cooperation is productive.

This is why companies exist even in a world of highly capable individuals. Coordination through repeated market contracts is not always the most efficient way to produce complex things. A stable organization can develop shared knowledge, specialize roles, invest in equipment, maintain standards, and make commitments that extend beyond 1 transaction. But organizations can also become slow, political, and wasteful. The founder must continually judge which activities should be done inside the firm, through partners, by independent specialists, or not at all.

\subsection{Incentive Alignment and Institutional Architecture}

Incentives matter because people do not automatically share the founder's priorities. Employees care about compensation, growth, meaningful work, fairness, and the predictability of their lives. Customers care about receiving what was promised. Investors care about risk and returns. Regulators and communities care about compliance and broader impacts. Good leadership attempts to align these interests where possible and manage unavoidable conflicts honestly.

Consider the difference between offering an employee a small share of equity and offering them a meaningful role in decisions. Equity may encourage identification with long-term results, but poorly structured ownership can be confusing or worthless. Authority without accountability produces chaos. Accountability without authority produces resentment. The real task is to design a system in which effort, judgment, information, and reward support one another.

\textbf{The founder's unusually valuable contribution may eventually be neither a technical solution nor an inspiring speech. It may be building an institution in which other people can do better work together than they could accomplish separately.}

This suggests a more ambitious conception of leverage: the compounding of other people's capabilities through a trusted organization.
